\section{Benchmark}
\label{sec:benchmark}

\subsection{Source datasets}
No public dataset depicts falls in industrial settings under a research-permitting licence, so we build the benchmark from public indoor fall datasets (Table~\ref{tab:datasets}). Two complement each other. UR Fall \cite{kwolek2014urfall} offers per-frame posture labels, a camera parallel to the floor and, for falls, a ceiling-mounted second camera that affords the overhead view common in industrial installations. The Multiple Cameras Fall dataset (MCF) \cite{auvinet2010mcfd} films each of 24 scenarios, performed by one actor, with eight elevated wide-angle cameras, and includes sitting, lying on a sofa and housekeeping that supply hard negatives; we use the dense re-annotation published with OmniFall \cite{schneider2025omnifall}. We also intended to use UP-Fall \cite{martinez2019upfall} and Le2i \cite{charfi2013le2i}, but UP-Fall's public links returned a quota error and Le2i's repository was unavailable throughout; both are deferred. The two datasets used show staged events by actors in homes, offices or a laboratory (Section~\ref{sec:limitations}).

\begin{table*}[t]
\centering
\caption{Source datasets. Licence terms as stated on the official distribution pages on 20~October~2026.}
\label{tab:datasets}
\footnotesize
\setlength{\tabcolsep}{4pt}
\begin{tabular}{llllll}
\toprule
Dataset & Cameras & Resolution / fps & Viewpoints & Labels used & Licence \\
\midrule
UR Fall \cite{kwolek2014urfall} & 2 Kinect & $640{\times}480$ / 30 & side; ceiling (falls only) & per-frame lying / not lying (depth-based) & CC BY-NC-SA 4.0 \\
MCF \cite{auvinet2010mcfd} & 8 IP, wide-angle & $720{\times}480$ / 30$^\ast$ & elevated, oblique & OmniFall segments \cite{schneider2025omnifall} & none stated \\
\bottomrule
\end{tabular}

\smallskip
\raggedright\footnotesize $^\ast$The AVI headers state 120\,fps; frame counts match OmniFall's timeline at 30\,fps, which we use.
\end{table*}

\subsection{Labelling policy}
Each sample receives one of three labels. A sample is \pd{} when a person's torso lies on the floor or ground, however it got there; this covers falls and deliberate lying on the floor. UR Fall counts deliberate lying as an activity of daily living, hence a negative; a single frame cannot distinguish the two, and for a safety alarm both warrant a check, so we label them \pd{} with the subtype \textit{lying\_floor\_deliberate} and report them separately. Junca et al.\ met the same ambiguity: a prompt that caught more falls also flagged deliberate lie-downs \cite{junca2026thermal}.

A sample is a \textit{hard negative} when a person adopts a low or horizontal posture with the torso off the floor: crouching, kneeling (including on hands and knees), bending over, sitting (chair, sofa, bed or floor), or lying on a sofa or bed. A sample is \textit{normal} when every visible person stands or walks. Labels describe the image, not the actor's activity, so a frame with one person standing and another kneeling beside a mattress is a kneeling hard negative. We exclude transitional frames (mid-fall, getting down or up), frames with no visible person, and frames whose posture is equivocal, such as a fallen person propped up or a torso draped over a table.

\subsection{Construction}
Candidates are proposed automatically and each is reviewed visually. For UR Fall they derive from the official per-frame depth labels and the ratio of current to standing height; the overhead camera, unlabelled, inherits the label of the side-camera frame with the same index (both streams have equal frame counts). For MCF, candidates are sampled inside OmniFall segments labelled \textit{fallen}, \textit{sitting}, \textit{lying}, \textit{standing}, \textit{walk} and \textit{other}, at least 0.25\,s from segment boundaries. Per source sequence we take at most three frames per label (MCF: two per scenario, camera and label, one for walking), at least 1\,s apart, then subsample MCF to a fixed quota per label.

An AI assistant (Claude, Anthropic) conducted the visual review under the author's direction, inspecting contact sheets of every candidate and recording an accept, reject or relabel decision, with a reason, for each frame, per the policy above. Because a vision--language model helped label a benchmark for vision--language models, we validated the labels twice. First, a blind check: \HumanCheckText{} Second, \iffullcheck the author went through every frame with its label displayed, confirming \FullConfirmed{}, changing \FullChanged{} and removing \FullRemoved{}. A visible label invites anchoring, so the blind check is the unbiased measure; the frames the author labelled differently in the blind check were all confirmed here, so the full pass is a screen for gross errors, not independent relabelling.\else\pending{every frame checked by eye by the author: confirmed / changed / removed counts}\fi{} Of \BenchNCand{} candidates, the review kept \BenchNKept{}, rejected \BenchNRejected{}, and relabelled \BenchNRelabelled{} of those kept. These corrections matter to anyone reusing the datasets: several UR Fall frames labelled ``lying on the ground'' show a person on hands and knees, ``upright'' frames include people bending over a sofa, and some MCF segments labelled \textit{sitting} show an empty view from cameras that cannot see the chair.

Every sample also carries a clip: the labelled frame plus the seven preceding frames at 0.25\,s spacing (1.75\,s in all), so decisions on 1, 4 or 8 frames end at the same labelled frame. Each sample records its viewpoint: \textit{side} (UR Fall camera~0), \textit{overhead} (UR Fall camera~1) or \textit{elevated} (MCF).

\subsection{Splits}
Frames from one sequence are strongly correlated, and MCF films one moment from eight cameras. We therefore split by sequence: an MCF scenario with all its cameras, or a UR Fall sequence with both, goes wholly to development or wholly to test. About a quarter of the sequences of each kind form the development set, which serves solely to choose thresholds. The test set has \BenchNTest{} samples from \BenchNSeqTest{} sequences (\BenchTestDown{} \pd{}, \BenchTestHard{} hard negative, \BenchTestNormal{} normal), the development set \BenchNDev{} samples from \BenchNSeqDev{} sequences (Table~\ref{tab:benchmark}). Because UR Fall's overhead camera recorded only falls, overhead frames are almost uniformly \pd{}: for that viewpoint we report recall but not a false-alarm rate.

\begin{table}[t]
\centering
\caption{Benchmark composition after visual review.}
\label{tab:benchmark}
\footnotesize
\setlength{\tabcolsep}{2.4pt}
\begin{tabular}{llrrrr}
\toprule
Label & Subtype & UR Fall & MCF & Test & Dev \\
\midrule
person down & fall & 82 & 136 & 160 & 58 \\
 & lying on floor, deliberate & 15 & 0 & 10 & 5 \\
\midrule
hard negative & crouching & 11 & 20 & 22 & 9 \\
 & kneeling (incl.\ hands and knees) & 10 & 11 & 16 & 5 \\
 & bending & 14 & 29 & 25 & 18 \\
 & sitting on a chair & 8 & 53 & 47 & 14 \\
 & sitting on a sofa or bed & 9 & 41 & 19 & 31 \\
 & sitting on the floor & 1 & 0 & 1 & 0 \\
 & lying on a sofa or bed & 18 & 38 & 41 & 15 \\
\midrule
normal & standing & 62 & 33 & 65 & 30 \\
 & walking & 0 & 120 & 88 & 32 \\
\midrule
Total & & 230 & 481 & 494 & 217 \\
\bottomrule
\end{tabular}

\end{table}

\subsection{Reproducibility}
UR Fall is licensed for non-commercial use with share-alike terms and MCF states no licence, so we do not redistribute frames. We release the label files (dataset, sequence, camera, frame index, label, subtype, viewpoint, review note) and scripts that download the source data from the official pages and rebuild every sample. The UR Fall script fetches only the needed frames through HTTP range requests, keeping the download below 1\,GB.
